\section{Choosing signs at the boundary}\label{stab:sign:section}

The preceding tail estimates control cancellation by comparing signed
and absolute means. We now retain more of the individual segment
representations. Two sign choices will be useful. The first assigns
one sign to every possible last vertex in the head; the second assigns
one sign to every first vertex outside it. Both produce tail vectors
whose norms can be recovered from their boundary coordinates.

The preferred-parent argument applies to either coordinate model $Z$.
For the subsequent exit-sign and trimming arguments we use the
infinite-height coordinate predual $B_\infty$ from
Section~\ref{binf:definition}.  Finite representations and head tests
come from the model's finite-hull theorem; the coordinate projections
are contractive and finite coordinate vectors are norm dense.
Write $H$ for a finite ancestral set. Its gates $G_H$ are the first
vertices outside $H$, including roots of untouched forest components.
A segment crossing the boundary has an exterior portion $[g,w]$ with
$g\in G_H$; the parent of $g$ is its last head vertex.

The shared coefficient estimate \eqref{stab:sign:coefficients}, rooted-mass
bound \eqref{stab:sign:absolute-masses} and same-sign norm formula
\eqref{stab:sign:rooted-equality} will turn crossing coefficients into
tail norms. We also use the two subtree estimates
\eqref{stab:cone-bounds}; the lower estimate reads the root coordinates,
and the upper estimate combines the norms in incomparable subtrees.

\subsection{One preferred sign at each last head vertex}

Grouping exits by their parent permits different child subtrees to
share one sign. The loss is controlled by the conflict between two
nested head segments ending at that parent.

\begin{proposition}\label{stab:sign:parent-lens}
Let $H$ be finite and ancestral, $x=P_Hx$ with $\|x\|=1$, and
$u\in\ker P_H$. For every $L>1$,
\begin{equation}\label{stab:sign:parent-lens-bound}
 \max\{\|x+u\|,\|x-u\|\}<L
 \quad\Longrightarrow\quad
 \|u\|\le
 3\left(\sqrt{L^2-1}
            +2\sqrt{\frac{1-L^{-2}}2}\right).
\end{equation}
\end{proposition}

\begin{proof}
First suppose $u$ has finite support. The finite-hull theorem for
$Z$ gives two finite laws of represented vectors
$A=\sum_i c_i d_{S_i}$ with disjoint segments and
$\sum_i c_i^2\le L^2$, such that
$\mathbb E_+A=x+u$ and $\mathbb E_-A=x-u$.
In $B_\infty$ these laws also follow directly from the strict gauge
representation, since both endpoint norms are strictly less than $L$.
Choose a head test $p=P_Hp$ with $\rho(p)=1$ and $x(p)=1$.
For each represented outcome set
\[
 s_i=p(S_i),\qquad D=1-\sum_i s_i^2.
\]
Every such disjoint family is a test for $\rho$, so $D\ge0$.
Both conditional expectations of $\sum_i c_is_i$ equal one.
Expansion and Cauchy--Schwarz give, separately for both laws,
\begin{equation}\label{stab:sign:parent-error}
 \mathbb E_\pm\sum_i(c_i-s_i)^2\le L^2-1,
 \qquad
 \mathbb E_\pm D\le1-L^{-2}.
\end{equation}
For the second inequality, use
$1\le L(\mathbb E_\pm\sum_i s_i^2)^{1/2}$.

Pool the laws with probabilities $1/2$, and write $\mathbb E$ for
their mixture. At a vertex $v\in H$ there is at most one segment of
an outcome that crosses with last head vertex $v$: two such segments would
both contain $v$. Let $C_v$ be its $s_i$ value, or zero if it is
absent. Define
\[
 p_v=\mathbb E(C_v)_+,\qquad n_v=\mathbb E(-C_v)_+.
\]
Take an independent second outcome, denoted by primes. If
$C_vC'_v<0$, their head portions are distinct nested segments ending
at $v$. Split the longer immediately before the shorter starts.
If its old sum is $C_v$, the new sums are $C_v-C'_v,C'_v$;
the gain in squared sums is
\[
 (C_v-C'_v)^2+(C'_v)^2-C_v^2
       =2(C'_v)^2-2C_vC'_v\ge2|C_vC'_v|.
\]
Make every such split in the outcome having the longer head.
Different $v$ require different original segments, so these splits
preserve disjointness in both outcomes. Their total gain is at most
$D+D'$. Independence and \eqref{stab:sign:parent-error} yield
\begin{equation}\label{stab:sign:parent-conflicts}
 2\sum_v p_vn_v\le1-L^{-2}.
\end{equation}

Choose sign $+$ at $v$ if $p_v\ge n_v$ and sign $-$ otherwise.
The minority masses $m_v=\min(p_v,n_v)$ satisfy
\begin{equation}\label{stab:sign:parent-minority}
 \sum_vm_v^2\le\frac{1-L^{-2}}2.
\end{equation}
In $Q_HA$, replace each coefficient $c_i$ by $s_i$. In the
expectation of either law this changes the vector by at most
$\sqrt{L^2-1}$, by \eqref{stab:sign:coefficients}, contractivity
of $Q_H$, and \eqref{stab:sign:parent-error}. Segments wholly outside
$H$ disappear because $s_i=0$. Split the remaining expected crossing
terms into $G_\pm+B_\pm$, retaining in $G_\pm$ the preferred
parent sign and putting the other terms in $B_\pm$.

For children $g\in G_H$ of a fixed $v$, the sum of absolute
coefficients in $B_\pm$ is at most $2m_v$: a nonnegative
conditional expectation is at most twice its mixture expectation.
Put $b_g=|B_\pm(e_g)|$. Formula~\eqref{stab:sign:rooted-equality},
followed by $\sum_g b_g^2\le(\sum_g b_g)^2$ within each parent
group, gives
\[
 \|B_\pm\|\le2\left(\sum_vm_v^2\right)^{1/2}
       \le2\sqrt{(1-L^{-2})/2}.
\]
Thus $\pm u=G_\pm+E_\pm$, where
\[
 \|E_\pm\|\le b_L,
 \qquad b_L=\sqrt{L^2-1}+2\sqrt{(1-L^{-2})/2}.
\]
The two $G$ vectors have the same sign at each exit vertex.
Lemma~\ref{stab:sign:rooted} therefore gives
$\|G_+\|\le\|G_++G_-\|\le2b_L$, and hence
$\|u\|\le3b_L$.

For a completed tail $u$, approximate it by $Q_Hu_j$ with finitely
supported $u_j$. The strict endpoint hypothesis persists for all
sufficiently large $j$. Passing to the limit proves the proposition.
\end{proof}

We use the averaged modulus $\widehat\delta_Z$ from
\eqref{eq:averaged-modulus}.

\begin{corollary}\label{stab:sign:parent-modulus}
Let $b_L=\sqrt{L^2-1}+2\sqrt{(1-L^{-2})/2}$ for $L>1$.
If $0<t<1$ and $L$ is chosen so that $3b_L<t$, then
\[
 \widehat\delta_Z(t)\ge (L-1)/8>0.
\]
\end{corollary}

\begin{proof}
Put $\delta=(L-1)/8$. Given a unit $x$, choose a finite-head unit
$x_0$ with $\|x-x_0\|<\delta$, supported on an ancestral $H$.
If $y\in\ker P_H$, $\|y\|\ge1$, and the average of
$\|x\pm ty\|$ were at most $1+\delta$, the average at $x_0$
would be less than $1+2\delta$. Each endpoint at $x_0$ has norm at
least one by $P_H$-contractivity. Their maximum would therefore be
less than $1+4\delta<L$. Proposition~\ref{stab:sign:parent-lens}
would give $t\le\|ty\|\le3b_L<t$, a contradiction.
\end{proof}

\begin{corollary}[A forest lens from preferred parent signs]
\label{stab:sign:forest-parent-lens}
In the forest space $X$, let $h=P_Hh$, $P_Hu=0$, and
$h+u,h-u\in B_X$.  Then
\begin{equation}\label{stab:sign:forest-parent-bound}
 \|u\|\le3(\sqrt2+2)\sqrt{1-\|h\|}.
\end{equation}
\end{corollary}

\begin{proof}
If $h=0$, then $\|u\|\le1$ proves the assertion.  Otherwise fix
a norming head test $p$, so $\rho(p)=1$ and $h(p)=\|h\|$.
For $\epsilon>0$, approximate $h\pm u$ within $\epsilon$ by
finite convex combinations $V_\pm=\mathbb E_\pm A$ of atoms.
Put $d_\epsilon=1-\|h\|+\epsilon$ and mix the two laws equally.
The individual evaluation deficits obey
$\mathbb E_\pm L(A)\le d_\epsilon$.  Therefore
Lemma~\ref{stab:coef:pointwise} gives the prelimit budgets
\begin{equation}\label{stab:sign:parent-approx-budget}
 \mathbb E_\pm E(A)\le2d_\epsilon,\qquad
 \mathbb E_\pm D(A)\le2d_\epsilon.
\end{equation}

Group only the crossings by their last head vertex $v$, and let
$p_v,n_v$ be their positive and negative mean head-sum masses in
the mixture.  The simultaneous parent-indexed splitting in the
proof of Proposition~\ref{stab:sign:parent-lens} gives
$\sum_vp_vn_v\le d_\epsilon$.  This step compares crossings
that may leave through different children of the same parent.
The proof applies in each forest component separately; untouched
roots have no crossing term.

Choose the same preferred parent sign for both laws.  Each law's
wrong-sign mass at $v$ is at most $2\min(p_v,n_v)$, so the sum of
its squared parent masses is at most $4d_\epsilon$.  Splitting each
parent mass among its child gates cannot increase that square sum.
Thus the discarded replaced tail has norm at most $2\sqrt{d_\epsilon}$
by \eqref{stab:sign:rooted-equality}.  Replacement itself costs at
most $\sqrt{2d_\epsilon}$ under each law.  If $M_\pm$ are the
retained replaced tails, contractivity and the approximation errors
give
\begin{equation}\label{stab:sign:parent-approx-remainder}
 \|\pm u-M_\pm\|
 \le\epsilon+(\sqrt2+2)\sqrt{d_\epsilon}.
\end{equation}
The signs of $M_+$ and $M_-$ agree at every gate.  Hence their
rooted norm formula gives
$\|M_+\|\le\|M_++M_-\|$, at most twice the right side of
\eqref{stab:sign:parent-approx-remainder}.  Adding the remaining
error and letting $\epsilon\downarrow0$ proves
\eqref{stab:sign:forest-parent-bound}.
\end{proof}

\subsection{Majority signs at the first exit}

From this point through the trimming argument we work in $B_\infty$,
and take the finite ancestral heads to contain its root.

There is a second way to remove the minority contributions. We keep
the original atom coefficients, assigning signs at individual exit
vertices. The cost now separates into a wholly exterior part and a
wrong-sign crossing part. This leads to a direct quadratic bound for
the averaged modulus.

\begin{proposition}\label{stab:sign:exit-lens}
Let $x=P_Hx$, $\|x\|=1$, and $y\in\ker P_H$. Put
$C_{\mathrm{ex}}=2+3\sqrt2$. For every $\delta>0$, with
$\eta=\delta/(1+\delta)$,
\begin{equation}\label{stab:sign:exit-lens-bound}
 \max\{\|x+y\|,\|x-y\|\}\le1+\delta
 \quad\Longrightarrow\quad
 \frac{\|y\|}{1+\delta}\le3C_{\mathrm{ex}}\sqrt\eta.
\end{equation}
\end{proposition}

\begin{proof}
First suppose $y$ is supported on $K\setminus H$ for a finite
ancestral $K\supset H$. By Lemma~\ref{binf:finite-hull}, represent
$(x\pm y)/(1+\delta)$ as finite expectations $\mathbb E_\pm A$
of atoms $A=\sum_i c_id_{S_i}$ supported in $K$, where
$\sum_i c_i^2\le1$. Choose a norming head test $p$ for $x$,
set $s_i=p(S_i)$, and put $D_A=1-\sum_i s_i^2$.
The shared identity \eqref{stab:deficit-identity} gives
\[
 D_A+\sum_i(c_i-s_i)^2\le2(1-A(p)).
\]
Both terms on the left are nonnegative. Since
$\mathbb E_\pm A(p)=1/(1+\delta)$, both conditional laws and
their equal mixture satisfy
\begin{equation}\label{stab:sign:exit-error}
 \mathbb E_\pm\sum_i(c_i-s_i)^2\le2\eta,
 \qquad \mathbb E_\pm D_A\le2\eta.
\end{equation}

For a crossing at $g\in G_H\cap K$, denote its coefficients by
$c_g=c_i$ and $s_g=s_i$, using zero for an absent crossing.
Take independent outcomes from the mixture. The full opposite-sign
comparison \eqref{stab:coef:head-conflict-eq}, applied to their head
sums, gives
\[
 2\sum_g(-s_gs'_g)_+\le D_A+D_{A'}.
\]
Its expectation is $4\sum_g\mathbb E(s_g)_+\,\mathbb E(-s_g)_+$
on the left and at most $4\eta$ on the right. Consequently
\begin{equation}\label{stab:sign:exit-conflicts}
 \sum_g\mathbb E(s_g)_+\,\mathbb E(-s_g)_+\le\eta.
\end{equation}
Choose $\sigma_g\in\{-1,1\}$ to be a majority sign for $s_g$.
The Euclidean norm of the vector of expected wrong-sign $s_g$
masses is at most $\sqrt\eta$, because
$\min(P,Q)^2\le PQ$. Moreover, Jensen's inequality in
Euclidean space and \eqref{stab:sign:exit-error} imply
\[
 \left\|(\mathbb E|c_g-s_g|)_g\right\|_2
       \le\sqrt{2\eta}.
\]
Using $(-\sigma_gc_g)_+\le(-\sigma_gs_g)_++|c_g-s_g|$,
we obtain
\begin{equation}\label{stab:sign:exit-minority}
 \left\|(\mathbb E(-\sigma_gc_g)_+)_g\right\|_2
       \le(1+\sqrt2)\sqrt\eta.
\end{equation}
For either conditional law the bound is twice this large.

We now recover the tail vectors from these masses. Wholly exterior
segments have $s_i=0$. Their expected vector in either law has norm
at most $\sqrt{2\eta}$ by \eqref{stab:sign:exit-error} and
\eqref{stab:sign:coefficients}. A crossing at $g$ has exterior
segment $L_g=[g,w]$. Its retained and discarded conditional means
are
\[
 g_{\pm,g}=\sigma_g\mathbb E_\pm
          [ (\sigma_g c_g)_+d_{L_g}],\qquad
 h_{\pm,g}=-\sigma_g\mathbb E_\pm
          [(-\sigma_g c_g)_+d_{L_g}],
\]
with zero terms for absent crossings. Since $\|d_{L_g}\|=1$,
\eqref{stab:cone-bounds} and
\eqref{stab:sign:exit-minority} bound
$\|\sum_gh_{\pm,g}\|$ by
$2(1+\sqrt2)\sqrt\eta$. Thus, with
$g_\pm=\sum_gg_{\pm,g}$,
\begin{equation}\label{stab:sign:exit-discard}
 \pm\frac{y}{1+\delta}=g_\pm+r_\pm,
 \qquad \|r_\pm\|\le C_{\mathrm{ex}}\sqrt\eta.
\end{equation}
To verify the required comparison explicitly, put
$C_{\pm,g}=\mathbb E_\pm(\sigma_gc_g)_+$. Then
$\|g_{\pm,g}\|\le C_{\pm,g}$ and
$g_{\pm,g}(e_g)=\sigma_gC_{\pm,g}$. The upper and lower
estimates in \eqref{stab:cone-bounds} give
\[
 \|g_+\|\le\|(C_{+,g})\|_2
 \le\|(C_{+,g}+C_{-,g})\|_2
 \le\|g_++g_-\|\le2C_{\mathrm{ex}}\sqrt\eta.
\]
Adding the remainder in \eqref{stab:sign:exit-discard} proves the
claim for finite support.

For a general completed tail, let $K$ increase through a finite
ancestral exhaustion containing $H$. Contractivity gives
$\|x\pm P_Ky\|\le1+\delta$, and $P_Ky\to y$ in norm by
coordinate-span density. The finite-support estimate passes to the
limit with the same $\delta$. This uses exact finite closed-ball
representations, rather than a strict-radius approximation.
\end{proof}

\begin{corollary}\label{stab:sign:exit-modulus}
For $0<t<1$ and $C_{\mathrm{ex}}=2+3\sqrt2$,
\[
 \widehat\delta_{B_\infty}(t)\ge\frac{t^2}{576C_{\mathrm{ex}}^2}.
\]
\end{corollary}

\begin{proof}
Set $\delta_t=(t/(12C_{\mathrm{ex}}))^2<1$. For a unit finite-head $x_0$
and $y\in\ker P_H$ with $\|y\|\ge1$, the two inequalities
$\|x_0\pm ty\|\le1+\delta_t$ would imply
\[
 \|ty\|\le3C_{\mathrm{ex}}\sqrt{\delta_t(1+\delta_t)}
       \le6C_{\mathrm{ex}}\sqrt{\delta_t}=t/2,
\]
which is impossible. Each endpoint has norm at least one, so their
average exceeds $1+\delta_t/2$. For an arbitrary unit $x$, choose
a finite-head unit $x_0$ with $\|x-x_0\|<\delta_t/4$.
The resulting average at $x$ exceeds $1+\delta_t/4$, uniformly
on this tail subspace. Formula~\eqref{eq:averaged-modulus}
gives the conclusion.
\end{proof}

\section{Trimming by maximal splitting gain}\label{stab:trim:section}

The next argument keeps the absolute crossing masses of the original
coefficients. A crossing is discarded if its coefficient is too
small, if it is poorly aligned with the head test, or if its head
allows a large increase by splitting. The remaining crossings have
one sign in each exit component. The point requiring care is the
discarded mass: a small probability alone does not control its
Euclidean norm.

\begin{proposition}\label{stab:trim:vanishing-lens}
There are no sequences of finite ancestral sets $H_n$, finite-head
unit vectors $x_n=P_{H_n}x_n$, and finitely supported
$z_n\in\ker P_{H_n}$ such that
\[
 \max\{\|x_n+z_n\|,\|x_n-z_n\|\}\longrightarrow1,
 \qquad \liminf_n\|z_n\|>0.
\]
The same assertion holds with $z_n$ in the completed tail spaces.
\end{proposition}

\begin{proof}
We first prove the finite-support assertion. Choose radii
$R_n>\max\{\|x_n+z_n\|,\|x_n-z_n\|\}$ with $R_n\to1$.
The gauge representation supplies two finite laws of atoms with
means $(x_n\pm z_n)/R_n$. Pool them equally. Suppress the index
$n$ and choose a norming head test $p$ for $x$.
For an outcome $A=\sum_i c_id_{S_i}$, put $s_i=p(S_i)$,
\[
 E_A=\sum_i(c_i-s_i)^2,\qquad D_A=1-\sum_i s_i^2.
\]
Since $\mathbb E_\pm A(p)=1/R$, the same calculation as in
\eqref{stab:sign:exit-error} gives
\begin{equation}\label{stab:trim:error}
 \mathbb E_\pm E_A\longrightarrow0,\qquad
 \mathbb E_\pm D_A\longrightarrow0.
\end{equation}
All expectations in what follows are finite, even though the tree
has infinitely many children.

For a crossing at $g$, let $c_g,s_g$ be its coefficient and head
sum. Its head is a suffix $[v,\operatorname{parent}(g)]$ of the
ancestral path. Define $d_g\ge0$ as the maximum gain in the sum of
squared $p$-sums obtainable by splitting this head at one place,
leaving the exterior portion attached to the later piece. Include
the option of no split, with gain zero. The head is finite, so this
maximum exists. Define all three quantities to be zero if there is
no crossing at $g$. Maximizing splits for different crossings act
on distinct disjoint segments. Therefore
\begin{equation}\label{stab:trim:gain-budget}
 \sum_gd_g\le D_A.
\end{equation}

Let $W$ be the part of $A$ consisting of wholly exterior segments.
These have $s_i=0$, so
\begin{equation}\label{stab:trim:exterior}
 \|\mathbb E_\pm W\|\le(\mathbb E_\pm E_A)^{1/2}=o(1).
\end{equation}
Also the Euclidean norm of the expected root-coordinate vector
$(\mathbb EW(e_g))_g$ is $o(1)$. Indeed, exit vertices are
incomparable, so a wholly exterior segment meets at most one of
them; disjointness then bounds the root-coordinate norm by the
Euclidean coefficient norm of the exterior segments.

Put $m_g=\mathbb E|c_g|$. We will show $\|(m_g)\|_2=o(1)$.
Indices with $m_g=0$ can be omitted. Discard an observation at $g$
if any one of the following holds:
\begin{equation}\label{stab:trim:tests}
 |c_g|<m_g/4,\qquad
 |c_g-s_g|>m_g/8,\qquad
 \sqrt{d_g}\ge m_g/8.
\end{equation}
The first condition removes an expected absolute mass at most
$m_g/4$. For the union $B_g$ of the last two conditions, put
$p_g=\mathbb P(B_g)$ and $b_g=\mathbb E[|c_g|1_{B_g}]$.
Markov's inequality, \eqref{stab:trim:error}, and
\eqref{stab:trim:gain-budget} imply
\begin{equation}\label{stab:trim:bad-probability}
 \sum_gm_g^2p_g
 \le64\mathbb E\sum_g\big((c_g-s_g)^2+d_g\big)=o(1).
\end{equation}
Write the last quantity as $\varepsilon_n\to0$.

To bound actual mass removal, fix $\eta>0$ and separate the
components according to $p_g>\eta$ or $p_g\le\eta$.
On the first set, $b_g\le m_g$, hence
\[
 \sum_{p_g>\eta}b_g^2\le\varepsilon_n/\eta.
\]
On the second set, Cauchy--Schwarz gives
$b_g^2\le p_g\mathbb E[c_g^21_{B_g}]$; therefore
\begin{equation}\label{stab:trim:bad-mass}
 \sum_{p_g\le\eta}b_g^2
       \le\eta\mathbb E\sum_gc_g^2\le\eta,
 \qquad
 \|b\|_2^2\le\varepsilon_n/\eta+\eta.
\end{equation}
Choose $\eta=\eta_n\to0$ slowly enough that
$\varepsilon_n/\eta_n\to0$ (for example
$\eta_n=\sqrt{\varepsilon_n}+1/n$). Thus $\|b\|_2=o(1)$.

Every surviving observation has $c_g$ and $s_g$ of the same sign,
$|s_g|\ge m_g/8$, and $d_g<(m_g/8)^2$. If two survivors at
one exit vertex had opposite signs, their head suffixes would be
distinct and nested. Splitting the longer at the shorter would
increase its squared sum by at least $2(m_g/8)^2$, contradicting
the definition and bound of its maximal gain. Hence the survivors
have one common sign at each $g$.

Let $a_g$ be the mass discarded by the first condition; thus
$a_g\le m_g/4$. The total discarded mass is at most $a_g+b_g$,
so the single-sign conclusion implies
\begin{equation}\label{stab:trim:signed-mean}
 |\mathbb Ec_g|\ge m_g-2(a_g+b_g)
                 \ge m_g/2-2b_g.
\end{equation}
The mixture has mean $x/R$ and therefore zero exterior coordinates.
At the exit roots its crossing part has values $\mathbb Ec_g$;
the remaining values come from $\mathbb EW$. By
\eqref{stab:trim:exterior}, $\|(\mathbb Ec_g)_g\|_2=o(1)$.
Equation~\eqref{stab:trim:signed-mean} consequently yields
\begin{equation}\label{stab:trim:mass-vanishes}
 \tfrac12\|m\|_2
       \le\|(\mathbb Ec_g)_g\|_2+2\|b\|_2=o(1).
\end{equation}

The expected absolute crossing mass at each $g$ in either
conditional law is at most $2m_g$. By
\eqref{stab:sign:absolute-masses}, its expected crossing tail has
norm at most $2\|m\|_2=o(1)$. Add the wholly exterior part from
\eqref{stab:trim:exterior}. Since the conditional tail means are
$\pm z/R$, we obtain $\|z/R\|=o(1)$, the desired contradiction.

For completed tails, choose finitely supported approximations and
apply $Q_{H_n}$, with errors tending to zero. Contractivity preserves
the tail condition and the endpoint limits, while the lower norm
bound persists. The finite-support contradiction applies.
\end{proof}

\begin{corollary}\label{stab:trim:modulus}
For the coordinate segment space, $\widehat\delta_{B_\infty}(t)>0$ for
every $0<t<1$.
\end{corollary}

\begin{proof}
The modulus is nonnegative by convexity. If it vanished at $t$,
there would be unit vectors $v_n$ for which the supremum over
finite-codimensional subspaces in
\eqref{eq:averaged-modulus} tends to zero. Choose unit
finite-head $x_n$ with $\|x_n-v_n\|<1/n$ and support contained
in an ancestral $H_n$. Using $\ker P_{H_n}$ in that supremum,
choose $z_n$ there with $\|z_n\|\ge t$ and
\[
 \frac{\|x_n+z_n\|+\|x_n-z_n\|}{2}\longrightarrow1.
\]
The choice of $z_n$ uses approximate infima; attainment is unnecessary.
Both endpoint norms are at least one by projection, and therefore
both converge to one. Approximation by finite tails gives, if desired,
$\|z_n\|\ge t-o(1)$ with the same endpoint limit. This contradicts
Proposition~\ref{stab:trim:vanishing-lens}.
\end{proof}
